\section{Cómo evaluar: pass@k, 10@k y el ranking real}

\subsection{La cobertura de la búsqueda y la selección del envío se deben ver por separado}

\videofigure{figures/pass-k-metrics.jpg}{pass@$k$ mide la cobertura de candidatos, 10@$k$ mide la tasa de éxito real tras elegir diez envíos a partir de $k$ muestras.}{00:14:18--00:16:26}

Si entre $n$ muestras hay $c$ programas correctos, al elegir $k$ sin reemplazo la estimación habitual es:

$$
\operatorname{pass@}k=1-\frac{\binom{n-c}{k}}{\binom{n}{k}}.
$$

\begin{itemize}
    \item $n$: número total de candidatos;
    \item $c$: de ellos, los candidatos que pasan las pruebas ocultas;
    \item $k$: el presupuesto de candidatos que se permite observar o extraer;
    \item pass@$k$: la probabilidad de que exista al menos un candidato correcto.
\end{itemize}

10@$k$ se acerca más a una competencia real: primero se generan $k$ candidatos y luego, mediante filtrado y agrupamiento (clustering), se eligen hasta diez envíos. Evalúa a la vez al generator y al selector, y no solo si «la respuesta correcta llegó a aparecer».

\videofigure{figures/codecontests-results.jpg}{Los resultados de CodeContests muestran que tanto el tamaño del modelo como el agrupamiento (clustering) mejoran el solve rate con un presupuesto de envíos limitado.}{00:16:26--00:18:08}

\begin{warningbox}{Un pass@$k$ alto no significa que el sistema sea desplegable}
Si el candidato correcto está escondido entre cientos de miles de programas pero el selector no puede identificarlo, no tiene ningún valor para el usuario real. Las métricas de despliegue deben incluir explícitamente el presupuesto de envíos, el costo de ejecución, la latencia y la tasa de selección errónea.
\end{warningbox}

\subsection{Evaluación real en Codeforces}

\videofigure{figures/codeforces-evaluation.jpg}{AlphaCode generó, filtró y envió programas en vivo en diez grandes competencias de Codeforces, con un ranking promedio cercano a la mediana de los participantes.}{00:09:04--00:13:58}

El curso enfatiza que este tipo de evaluación es más confiable que un benchmark estático: los problemas aparecieron después de que el modelo terminó su entrenamiento, el sistema opera con las reglas reales de envío y se compara contra miles de participantes humanos. AlphaCode alcanzó en su momento un ranking aproximado del 54\% superior, lo que demostró que la generación de programas de extremo a extremo ya no es solo una tarea de juguete.

\videofigure{figures/alphacode-tradeoffs.jpg}{Las fortalezas de AlphaCode son la resolución completa de problemas y la verificación reproducible; sus debilidades son el alto costo computacional, la dependencia de pruebas ocultas y el cuello de botella de la selección.}{00:22:42--00:24:34}

\subsection{Resumen de la sección}

Al evaluar sistemas de búsqueda hay que separar tres cosas: si el modelo puede generar un candidato correcto, si el selector puede encontrarlo con un presupuesto reducido y si el programa final puede pasar en una plataforma real con una eficiencia aceptable.

